\scanpage{7}
\mspage{5}

Pour ce dernier point, on est ramené au

\uline{Lemme} $A'$ fid plat sur $A$, $B$ une $A$-algèbre. Alors $B$ fini (entier) sur $A$ sss $B'$ l'est sur $A'$.

Fini est une histoire de module. Entier s'y ramène en regardant les $A[x]$, $x\in B$.

\textcircled{4} \uline{Descente de l'amplitude}

\diagram{\begin{tikzcd}[ampersand replacement=\&]
\underline L\quad X\arrow[d,dash] \& X'\arrow[d,dash]\\
Y\arrow[r,dash]\arrow[d,dash] \& \unsure{Y''|Y'}\\
S \& S'\arrow[l,"{\substack{f\ \text{fid plat}\\\text{qu cpt}}}"']
\end{tikzcd}}

\uline{$L$} (très) ample rel /$Y$ $\Longleftrightarrow$ $L$ \inserted{(très)} ample rel. à $Y$

On peut supposer $X$ cohérent sur $Y$, alors $g_*(L)$ \unsure{qu. cpt|qu. coh} et $g_*(L)'\simeq g'_*(L')$

Très ample signifie que 1) $g^*(g_*(L))\longrightarrow L$ épi, donc $X\mathrel{-}P(g_*(L))$ défini. \deleted{\unsure{2)}} 2) \deleted{$g^*(g_*(L))\unsure{\mathrel{-}L}$} $X\mathrel{-}P(g_*(L))$ une imm.

Si $g$ de t.f., ``ample'' se ramène : ``très ample'', donc c'est.

Pour le cas général, \inserted{on peut supposer $Y$ qu. cpt,} on utilise le critère \unsure{que|avec} pour $F$ qu. ch. de t.f. sur $S$. et $n$ grand, $f^*(f_*(F(n)))\mathrel{-}F(n)$ \unsure{épij.|épi}

On retrouve que $g$ quasi-affine $\Longleftrightarrow g'$ qu.-affine

\uline{NB} les seuls exemples de propriétés de morphismes qui ne se descendent pas par fpqc sont les morphismes \uline{projectifs} et \uline{quasi-projectifs}

\textcircled{5} \uline{Questions} d'\uline{effectivité} \uline{des données de descente}

% archive sheet number: 6
